% Options for packages loaded elsewhere
% Options for packages loaded elsewhere
\PassOptionsToPackage{unicode}{hyperref}
\PassOptionsToPackage{hyphens}{url}
\PassOptionsToPackage{dvipsnames,svgnames,x11names}{xcolor}
%
\documentclass[
]{report}
\usepackage{xcolor}
\usepackage[top=30mm,left=20mm,right=20mm,bottom=30mm]{geometry}
\usepackage{amsmath,amssymb}
\setcounter{secnumdepth}{5}
\usepackage{iftex}
\ifPDFTeX
  \usepackage[T1]{fontenc}
  \usepackage[utf8]{inputenc}
  \usepackage{textcomp} % provide euro and other symbols
\else % if luatex or xetex
  \usepackage{unicode-math} % this also loads fontspec
  \defaultfontfeatures{Scale=MatchLowercase}
  \defaultfontfeatures[\rmfamily]{Ligatures=TeX,Scale=1}
\fi
\usepackage{lmodern}
\ifPDFTeX\else
  % xetex/luatex font selection
\fi
% Use upquote if available, for straight quotes in verbatim environments
\IfFileExists{upquote.sty}{\usepackage{upquote}}{}
\IfFileExists{microtype.sty}{% use microtype if available
  \usepackage[]{microtype}
  \UseMicrotypeSet[protrusion]{basicmath} % disable protrusion for tt fonts
}{}
\makeatletter
\@ifundefined{KOMAClassName}{% if non-KOMA class
  \IfFileExists{parskip.sty}{%
    \usepackage{parskip}
  }{% else
    \setlength{\parindent}{0pt}
    \setlength{\parskip}{6pt plus 2pt minus 1pt}}
}{% if KOMA class
  \KOMAoptions{parskip=half}}
\makeatother
% Make \paragraph and \subparagraph free-standing
\makeatletter
\ifx\paragraph\undefined\else
  \let\oldparagraph\paragraph
  \renewcommand{\paragraph}{
    \@ifstar
      \xxxParagraphStar
      \xxxParagraphNoStar
  }
  \newcommand{\xxxParagraphStar}[1]{\oldparagraph*{#1}\mbox{}}
  \newcommand{\xxxParagraphNoStar}[1]{\oldparagraph{#1}\mbox{}}
\fi
\ifx\subparagraph\undefined\else
  \let\oldsubparagraph\subparagraph
  \renewcommand{\subparagraph}{
    \@ifstar
      \xxxSubParagraphStar
      \xxxSubParagraphNoStar
  }
  \newcommand{\xxxSubParagraphStar}[1]{\oldsubparagraph*{#1}\mbox{}}
  \newcommand{\xxxSubParagraphNoStar}[1]{\oldsubparagraph{#1}\mbox{}}
\fi
\makeatother


\usepackage{longtable,booktabs,array}
\usepackage{calc} % for calculating minipage widths
% Correct order of tables after \paragraph or \subparagraph
\usepackage{etoolbox}
\makeatletter
\patchcmd\longtable{\par}{\if@noskipsec\mbox{}\fi\par}{}{}
\makeatother
% Allow footnotes in longtable head/foot
\IfFileExists{footnotehyper.sty}{\usepackage{footnotehyper}}{\usepackage{footnote}}
\makesavenoteenv{longtable}
\usepackage{graphicx}
\makeatletter
\newsavebox\pandoc@box
\newcommand*\pandocbounded[1]{% scales image to fit in text height/width
  \sbox\pandoc@box{#1}%
  \Gscale@div\@tempa{\textheight}{\dimexpr\ht\pandoc@box+\dp\pandoc@box\relax}%
  \Gscale@div\@tempb{\linewidth}{\wd\pandoc@box}%
  \ifdim\@tempb\p@<\@tempa\p@\let\@tempa\@tempb\fi% select the smaller of both
  \ifdim\@tempa\p@<\p@\scalebox{\@tempa}{\usebox\pandoc@box}%
  \else\usebox{\pandoc@box}%
  \fi%
}
% Set default figure placement to htbp
\def\fps@figure{htbp}
\makeatother





\setlength{\emergencystretch}{3em} % prevent overfull lines

\providecommand{\tightlist}{%
  \setlength{\itemsep}{0pt}\setlength{\parskip}{0pt}}



 


\makeatletter
\@ifpackageloaded{tcolorbox}{}{\usepackage[skins,breakable]{tcolorbox}}
\@ifpackageloaded{fontawesome5}{}{\usepackage{fontawesome5}}
\definecolor{quarto-callout-color}{HTML}{909090}
\definecolor{quarto-callout-note-color}{HTML}{0758E5}
\definecolor{quarto-callout-important-color}{HTML}{CC1914}
\definecolor{quarto-callout-warning-color}{HTML}{EB9113}
\definecolor{quarto-callout-tip-color}{HTML}{00A047}
\definecolor{quarto-callout-caution-color}{HTML}{FC5300}
\definecolor{quarto-callout-color-frame}{HTML}{acacac}
\definecolor{quarto-callout-note-color-frame}{HTML}{4582ec}
\definecolor{quarto-callout-important-color-frame}{HTML}{d9534f}
\definecolor{quarto-callout-warning-color-frame}{HTML}{f0ad4e}
\definecolor{quarto-callout-tip-color-frame}{HTML}{02b875}
\definecolor{quarto-callout-caution-color-frame}{HTML}{fd7e14}
\makeatother
\makeatletter
\@ifpackageloaded{caption}{}{\usepackage{caption}}
\AtBeginDocument{%
\ifdefined\contentsname
  \renewcommand*\contentsname{Table of contents}
\else
  \newcommand\contentsname{Table of contents}
\fi
\ifdefined\listfigurename
  \renewcommand*\listfigurename{List of Figures}
\else
  \newcommand\listfigurename{List of Figures}
\fi
\ifdefined\listtablename
  \renewcommand*\listtablename{List of Tables}
\else
  \newcommand\listtablename{List of Tables}
\fi
\ifdefined\figurename
  \renewcommand*\figurename{Figure}
\else
  \newcommand\figurename{Figure}
\fi
\ifdefined\tablename
  \renewcommand*\tablename{Table}
\else
  \newcommand\tablename{Table}
\fi
}
\@ifpackageloaded{float}{}{\usepackage{float}}
\floatstyle{ruled}
\@ifundefined{c@chapter}{\newfloat{codelisting}{h}{lop}}{\newfloat{codelisting}{h}{lop}[chapter]}
\floatname{codelisting}{Listing}
\newcommand*\listoflistings{\listof{codelisting}{List of Listings}}
\makeatother
\makeatletter
\makeatother
\makeatletter
\@ifpackageloaded{caption}{}{\usepackage{caption}}
\@ifpackageloaded{subcaption}{}{\usepackage{subcaption}}
\makeatother
\usepackage{bookmark}
\IfFileExists{xurl.sty}{\usepackage{xurl}}{} % add URL line breaks if available
\urlstyle{same}
\hypersetup{
  pdftitle={Research Readiness Dashboard},
  colorlinks=true,
  linkcolor={blue},
  filecolor={Maroon},
  citecolor={Blue},
  urlcolor={Blue},
  pdfcreator={LaTeX via pandoc}}


\title{Research Readiness Dashboard}
\author{}
\date{}
\begin{document}
\maketitle

\renewcommand*\contentsname{Table of contents}
{
\hypersetup{linkcolor=}
\setcounter{tocdepth}{2}
\tableofcontents
}

\chapter{Purpose}\label{purpose}

Use this dashboard as a self-auditing summary to assess the structural
completeness, governance readiness, AI disclosure status, and
reproducibility of your research project.

\begin{center}\rule{0.5\linewidth}{0.5pt}\end{center}

\chapter{1. Project Overview Status}\label{project-overview-status}

\begin{longtable}[]{@{}
  >{\raggedright\arraybackslash}p{(\linewidth - 4\tabcolsep) * \real{0.3333}}
  >{\raggedright\arraybackslash}p{(\linewidth - 4\tabcolsep) * \real{0.3333}}
  >{\raggedright\arraybackslash}p{(\linewidth - 4\tabcolsep) * \real{0.3333}}@{}}
\toprule\noalign{}
\begin{minipage}[b]{\linewidth}\raggedright
Component
\end{minipage} & \begin{minipage}[b]{\linewidth}\raggedright
Status
\end{minipage} & \begin{minipage}[b]{\linewidth}\raggedright
Notes
\end{minipage} \\
\midrule\noalign{}
\endhead
\bottomrule\noalign{}
\endlastfoot
\textbf{Research question clearly defined} & ☑ & Outcome defined as
binary missed appointment. \\
\textbf{Dataset described} & ☑ & Synthetic dataset fully documented in
Data section. \\
\textbf{Methods documented} & ☑ & Logistic regression model specified
with equation. \\
\textbf{Results interpreted} & ☑ & Class imbalance and interpretation
discussed. \\
\textbf{Limitations stated} & ☑ & Limitations of synthetic data
acknowledged. \\
\end{longtable}

\begin{center}\rule{0.5\linewidth}{0.5pt}\end{center}

\chapter{2. Governance Completion}\label{governance-completion}

\begin{tcolorbox}[enhanced jigsaw, opacitybacktitle=0.6, coltitle=black, colback=white, opacityback=0, leftrule=.75mm, colbacktitle=quarto-callout-note-color!10!white, colframe=quarto-callout-note-color-frame, left=2mm, toprule=.15mm, breakable, titlerule=0mm, bottomtitle=1mm, rightrule=.15mm, bottomrule=.15mm, title=\textcolor{quarto-callout-note-color}{\faInfo}\hspace{0.5em}{Note}, toptitle=1mm, arc=.35mm]

These artefacts strengthen research transparency and defensibility.

\end{tcolorbox}

\begin{itemize}
\tightlist
\item[$\boxtimes$]
  Dataset Card completed
\item[$\boxtimes$]
  Model Card completed
\item[$\boxtimes$]
  Risk Register filled
\item[$\boxtimes$]
  Decision Log updated
\item[$\boxtimes$]
  Reproducibility Checklist completed
\end{itemize}

\emph{Note: All governance artefacts are available in the
\texttt{governance/} directory.}

\begin{center}\rule{0.5\linewidth}{0.5pt}\end{center}

\chapter{3. AI Disclosure \& Reflection}\label{ai-disclosure-reflection}

\begin{tcolorbox}[enhanced jigsaw, opacitybacktitle=0.6, coltitle=black, colback=white, opacityback=0, leftrule=.75mm, colbacktitle=quarto-callout-note-color!10!white, colframe=quarto-callout-note-color-frame, left=2mm, toprule=.15mm, breakable, titlerule=0mm, bottomtitle=1mm, rightrule=.15mm, bottomrule=.15mm, title=\textcolor{quarto-callout-note-color}{\faInfo}\hspace{0.5em}{Note}, toptitle=1mm, arc=.35mm]

AI use is \textbf{optional}. Whether used or not, human responsibility
remains primary. Transparency ensures rigorous research standards.

\end{tcolorbox}

\begin{itemize}
\tightlist
\item[$\boxtimes$]
  AI tools used? (Yes)
\item[$\boxtimes$]
  AI Capability Checkpoints completed
\item[$\boxtimes$]
  AI use documented in Decision Log
\item[$\boxtimes$]
  No sensitive data shared with AI tools
\item[$\boxtimes$]
  Human verification explicitly documented
\end{itemize}

\emph{Note: AI (a Large Language Model) was used for drafting assistance
only, specifically to write the Python synthetic data script and
critique the protocol. All outputs were reviewed and validated by the
human research team. No real patient data was involved.}

\begin{center}\rule{0.5\linewidth}{0.5pt}\end{center}

\chapter{4. Reproducibility Readiness}\label{reproducibility-readiness}

\begin{tcolorbox}[enhanced jigsaw, opacitybacktitle=0.6, coltitle=black, colback=white, opacityback=0, leftrule=.75mm, colbacktitle=quarto-callout-note-color!10!white, colframe=quarto-callout-note-color-frame, left=2mm, toprule=.15mm, breakable, titlerule=0mm, bottomtitle=1mm, rightrule=.15mm, bottomrule=.15mm, title=\textcolor{quarto-callout-note-color}{\faInfo}\hspace{0.5em}{Note}, toptitle=1mm, arc=.35mm]

A reproducible project ensures portability, transparency, and trust in
the scientific method.

\end{tcolorbox}

\begin{itemize}
\tightlist
\item[$\boxtimes$]
  Data dictionary completed
\item[$\boxtimes$]
  Codebook available
\item[$\boxtimes$]
  Synthetic vs real data clearly stated
\item[$\boxtimes$]
  Dependencies documented
\item[$\boxtimes$]
  Version control in use
\item[$\boxtimes$]
  Citation file included
\end{itemize}

\emph{Note: The project leverages Quarto and tracks dependencies
locally. Synthetic status is explicitly stated in the README and
limitations.}

\begin{center}\rule{0.5\linewidth}{0.5pt}\end{center}

\chapter{5. Summary Readiness
Indicator}\label{summary-readiness-indicator}

Use the following visual guide to interpret your total incomplete items
across all checklists above.

\begin{quote}
\textbf{0 incomplete items} → Research Ready for Internal Review

\textbf{1--5 incomplete items} → Requires Governance Review

\textbf{\textgreater5 incomplete items} → Early Draft Stage
\end{quote}

\textbf{This synthetic project demonstrates full structural, governance,
and AI transparency readiness within the template architecture.}




\end{document}
